\section{总结与延伸}

\subsection{核心洞察}

本次青稞RL专题讨论呈现了大模型RL领域当前最前沿的思考和实践经验。以下是贯穿全场讨论的核心洞察：

\begin{importantbox}{大模型RL的五大核心认知}
\begin{enumerate}
    \item \textbf{Infra决定上限}：2025年最具影响力的进展来自Infra而非算法。训推不一致、MoE稳定性等问题的解决方案都源于对底层系统的深刻理解
    \item \textbf{算法是Recipe}：单一算法的时代已过，取而代之的是一套经过验证的Trick组合。好的Recipe应该简单、符合直觉且有效
    \item \textbf{第一性原理不过时}：从GSPO的Sequence Level优化目标推导，到训推不一致的理论分析，回归基本原理始终是解决新问题的有力武器
    \item \textbf{经典RL理论的"文艺复兴"}：PPO中的Batch Norm、Policy Gradient的Variance Reduction、Off-Policy Correction等数十年前的技术在大模型场景下焕发新生
    \item \textbf{Co-design是趋势}：算法与Infra的联合优化、模型架构与RL训练的协同设计将成为常态
\end{enumerate}
\end{importantbox}

\subsection{开放问题}

讨论中暴露出若干尚未解决的开放问题：

\begin{enumerate}
    \item \textbf{Token-Level Value Model的根本困难}：从统计学习理论角度，训好Token-Level Value需要的采样量与序列长度呈多项式关系，在当前计算预算下存在Fundamental Limit
    \item \textbf{Exploration的缺失}：当前大模型RL几乎没有真正的Exploration，主要依赖模型自身的先验——未来如何引入结构化的探索机制？
    \item \textbf{Diffusion LLM对RL范式的影响}：如果Autoregressive范式被取代，RL的整个训练Pipeline可能需要重构
    \item \textbf{全异步训练下的Off-Policy容忍度}：给定一套算法配方，Off-Policy程度可以容忍到什么地步？
    \item \textbf{长程Agentic场景的Credit Assignment}：Agent执行数天任务时，如何有效地将最终Reward回传到关键决策点？
\end{enumerate}

\subsection{拓展阅读}

\begin{itemize}
    \item \textbf{PRIME}：崔干曲等人的工作，通过Implicit PRM实现自动Credit Assignment
    \item \textbf{Reinforce++}：胡建提出，回归PPO经典技巧的精简RL算法
    \item \textbf{GSPO}：郑楚杰等人提出的Sequence-Level优化算法，去掉KL Penalty提升训练效率
    \item \textbf{OpenRLHF}：开源RL训练框架，首创推理-训练引擎解耦架构（\url{https://github.com/OpenRLHF/OpenRLHF}）
    \item \textbf{TRPO}（Schulman et al., 2015）：Trust Region Policy Optimization，现代RL算法的理论基石
    \item \textbf{DAPO}：提出Clip Higher等Entropy控制方法的工作
    \item \textbf{Train-Inference Mismatch Blog}：李英儒团队2025年9月发布的关于训推不一致问题的深度分析
\end{itemize}

%% --- Fin del contenido --- %%

\end{document}
